% Options for packages loaded elsewhere
\PassOptionsToPackage{unicode}{hyperref}
\PassOptionsToPackage{hyphens}{url}
\documentclass[
]{article}
\usepackage{xcolor}
\usepackage[margin=1in]{geometry}
\usepackage{amsmath,amssymb}
\setcounter{secnumdepth}{-\maxdimen} % remove section numbering
\usepackage{iftex}
\ifPDFTeX
  \usepackage[T1]{fontenc}
  \usepackage[utf8]{inputenc}
  \usepackage{textcomp} % provide euro and other symbols
\else % if luatex or xetex
  \usepackage{unicode-math} % this also loads fontspec
  \defaultfontfeatures{Scale=MatchLowercase}
  \defaultfontfeatures[\rmfamily]{Ligatures=TeX,Scale=1}
\fi
\usepackage{lmodern}
\ifPDFTeX\else
  % xetex/luatex font selection
\fi
% Use upquote if available, for straight quotes in verbatim environments
\IfFileExists{upquote.sty}{\usepackage{upquote}}{}
\IfFileExists{microtype.sty}{% use microtype if available
  \usepackage[]{microtype}
  \UseMicrotypeSet[protrusion]{basicmath} % disable protrusion for tt fonts
}{}
\makeatletter
\@ifundefined{KOMAClassName}{% if non-KOMA class
  \IfFileExists{parskip.sty}{%
    \usepackage{parskip}
  }{% else
    \setlength{\parindent}{0pt}
    \setlength{\parskip}{6pt plus 2pt minus 1pt}}
}{% if KOMA class
  \KOMAoptions{parskip=half}}
\makeatother
\usepackage{color}
\usepackage{fancyvrb}
\newcommand{\VerbBar}{|}
\newcommand{\VERB}{\Verb[commandchars=\\\{\}]}
\DefineVerbatimEnvironment{Highlighting}{Verbatim}{commandchars=\\\{\}}
% Add ',fontsize=\small' for more characters per line
\usepackage{framed}
\definecolor{shadecolor}{RGB}{248,248,248}
\newenvironment{Shaded}{\begin{snugshade}}{\end{snugshade}}
\newcommand{\AlertTok}[1]{\textcolor[rgb]{0.94,0.16,0.16}{#1}}
\newcommand{\AnnotationTok}[1]{\textcolor[rgb]{0.56,0.35,0.01}{\textbf{\textit{#1}}}}
\newcommand{\AttributeTok}[1]{\textcolor[rgb]{0.13,0.29,0.53}{#1}}
\newcommand{\BaseNTok}[1]{\textcolor[rgb]{0.00,0.00,0.81}{#1}}
\newcommand{\BuiltInTok}[1]{#1}
\newcommand{\CharTok}[1]{\textcolor[rgb]{0.31,0.60,0.02}{#1}}
\newcommand{\CommentTok}[1]{\textcolor[rgb]{0.56,0.35,0.01}{\textit{#1}}}
\newcommand{\CommentVarTok}[1]{\textcolor[rgb]{0.56,0.35,0.01}{\textbf{\textit{#1}}}}
\newcommand{\ConstantTok}[1]{\textcolor[rgb]{0.56,0.35,0.01}{#1}}
\newcommand{\ControlFlowTok}[1]{\textcolor[rgb]{0.13,0.29,0.53}{\textbf{#1}}}
\newcommand{\DataTypeTok}[1]{\textcolor[rgb]{0.13,0.29,0.53}{#1}}
\newcommand{\DecValTok}[1]{\textcolor[rgb]{0.00,0.00,0.81}{#1}}
\newcommand{\DocumentationTok}[1]{\textcolor[rgb]{0.56,0.35,0.01}{\textbf{\textit{#1}}}}
\newcommand{\ErrorTok}[1]{\textcolor[rgb]{0.64,0.00,0.00}{\textbf{#1}}}
\newcommand{\ExtensionTok}[1]{#1}
\newcommand{\FloatTok}[1]{\textcolor[rgb]{0.00,0.00,0.81}{#1}}
\newcommand{\FunctionTok}[1]{\textcolor[rgb]{0.13,0.29,0.53}{\textbf{#1}}}
\newcommand{\ImportTok}[1]{#1}
\newcommand{\InformationTok}[1]{\textcolor[rgb]{0.56,0.35,0.01}{\textbf{\textit{#1}}}}
\newcommand{\KeywordTok}[1]{\textcolor[rgb]{0.13,0.29,0.53}{\textbf{#1}}}
\newcommand{\NormalTok}[1]{#1}
\newcommand{\OperatorTok}[1]{\textcolor[rgb]{0.81,0.36,0.00}{\textbf{#1}}}
\newcommand{\OtherTok}[1]{\textcolor[rgb]{0.56,0.35,0.01}{#1}}
\newcommand{\PreprocessorTok}[1]{\textcolor[rgb]{0.56,0.35,0.01}{\textit{#1}}}
\newcommand{\RegionMarkerTok}[1]{#1}
\newcommand{\SpecialCharTok}[1]{\textcolor[rgb]{0.81,0.36,0.00}{\textbf{#1}}}
\newcommand{\SpecialStringTok}[1]{\textcolor[rgb]{0.31,0.60,0.02}{#1}}
\newcommand{\StringTok}[1]{\textcolor[rgb]{0.31,0.60,0.02}{#1}}
\newcommand{\VariableTok}[1]{\textcolor[rgb]{0.00,0.00,0.00}{#1}}
\newcommand{\VerbatimStringTok}[1]{\textcolor[rgb]{0.31,0.60,0.02}{#1}}
\newcommand{\WarningTok}[1]{\textcolor[rgb]{0.56,0.35,0.01}{\textbf{\textit{#1}}}}
\usepackage{graphicx}
\makeatletter
\newsavebox\pandoc@box
\newcommand*\pandocbounded[1]{% scales image to fit in text height/width
  \sbox\pandoc@box{#1}%
  \Gscale@div\@tempa{\textheight}{\dimexpr\ht\pandoc@box+\dp\pandoc@box\relax}%
  \Gscale@div\@tempb{\linewidth}{\wd\pandoc@box}%
  \ifdim\@tempb\p@<\@tempa\p@\let\@tempa\@tempb\fi% select the smaller of both
  \ifdim\@tempa\p@<\p@\scalebox{\@tempa}{\usebox\pandoc@box}%
  \else\usebox{\pandoc@box}%
  \fi%
}
% Set default figure placement to htbp
\def\fps@figure{htbp}
\makeatother
\setlength{\emergencystretch}{3em} % prevent overfull lines
\providecommand{\tightlist}{%
  \setlength{\itemsep}{0pt}\setlength{\parskip}{0pt}}
\usepackage{bookmark}
\IfFileExists{xurl.sty}{\usepackage{xurl}}{} % add URL line breaks if available
\urlstyle{same}
\hypersetup{
  pdftitle={Visualisations et Analyse des Risques},
  hidelinks,
  pdfcreator={LaTeX via pandoc}}

\title{Visualisations et Analyse des Risques}
\author{}
\date{\vspace{-2.5em}}

\begin{document}
\maketitle

\section{1. Analyse de la Sinistralité
Temporelle}\label{analyse-de-la-sinistralituxe9-temporelle}

Pour bien comprendre la dynamique de notre portefeuille, la première
étape consiste à observer l'évolution du nombre de sinistres au fil des
années. Nous utilisons \texttt{ggplot2} pour visualiser cette
distribution.

\begin{Shaded}
\begin{Highlighting}[]
\FunctionTok{library}\NormalTok{(dplyr)}
\FunctionTok{library}\NormalTok{(ggplot2)}
\FunctionTok{library}\NormalTok{(lubridate)}

\CommentTok{\#1. chargement de la base propre fusionnée}

\NormalTok{base\_complete }\OtherTok{\textless{}{-}} \FunctionTok{readRDS}\NormalTok{(}\StringTok{"base\_complete\_clean.rds"}\NormalTok{)}
\end{Highlighting}
\end{Shaded}

\begin{Shaded}
\begin{Highlighting}[]
\NormalTok{graph\_annee }\OtherTok{\textless{}{-}}\NormalTok{ base\_complete }\SpecialCharTok{\%\textgreater{}\%}
  \CommentTok{\# On extrait l\textquotesingle{}année de la date de survenance du sinistre}
  \FunctionTok{mutate}\NormalTok{(}\AttributeTok{annee\_sinistre =} \FunctionTok{year}\NormalTok{(surv\_sin)) }\SpecialCharTok{\%\textgreater{}\%}
  \CommentTok{\# On filtre les éventuels NA sur la date}
  \FunctionTok{filter}\NormalTok{(}\SpecialCharTok{!}\FunctionTok{is.na}\NormalTok{(annee\_sinistre)) }\SpecialCharTok{\%\textgreater{}\%}
  \CommentTok{\# On compte le nombre de sinistres par année}
  \FunctionTok{count}\NormalTok{(annee\_sinistre, }\AttributeTok{name =} \StringTok{"total\_sinistres"}\NormalTok{) }\SpecialCharTok{\%\textgreater{}\%}
  
  \CommentTok{\# 3. Création du graphique avec ggplot2}
  \FunctionTok{ggplot}\NormalTok{(}\FunctionTok{aes}\NormalTok{(}\AttributeTok{x =} \FunctionTok{as.factor}\NormalTok{(annee\_sinistre), }\AttributeTok{y =}\NormalTok{ total\_sinistres)) }\SpecialCharTok{+}
  \FunctionTok{geom\_bar}\NormalTok{(}\AttributeTok{stat =} \StringTok{"identity"}\NormalTok{, }\AttributeTok{fill =} \StringTok{"\#2c3e50"}\NormalTok{) }\SpecialCharTok{+} \CommentTok{\# Un joli bleu professionnel}
  \FunctionTok{geom\_text}\NormalTok{(}\FunctionTok{aes}\NormalTok{(}\AttributeTok{label =}\NormalTok{ total\_sinistres), }\AttributeTok{vjust =} \SpecialCharTok{{-}}\FloatTok{0.5}\NormalTok{, }\AttributeTok{color =} \StringTok{"black"}\NormalTok{, }\AttributeTok{size =} \FloatTok{3.5}\NormalTok{) }\SpecialCharTok{+} \CommentTok{\# Ajout des étiquettes de données}
  \FunctionTok{labs}\NormalTok{(}
    \AttributeTok{title =} \StringTok{"Évolution du nombre de sinistres par année"}\NormalTok{,}
    \AttributeTok{subtitle =} \StringTok{"Base de données ENSAssuRances"}\NormalTok{,}
    \AttributeTok{x =} \StringTok{"Année de survenance"}\NormalTok{,}
    \AttributeTok{y =} \StringTok{"Total des sinistres"}
\NormalTok{  ) }\SpecialCharTok{+}
  \FunctionTok{theme\_minimal}\NormalTok{()}

\CommentTok{\# Affichage du graphique}
\FunctionTok{print}\NormalTok{(graph\_annee)}
\end{Highlighting}
\end{Shaded}

\pandocbounded{\includegraphics[keepaspectratio]{03_Visualisations_files/figure-latex/unnamed-chunk-1-1.pdf}}

\subsection{Interprétation : Évolution de la
sinistralité}\label{interpruxe9tation-uxe9volution-de-la-sinistralituxe9}

L'histogramme révèle des variations temporelles très marquées,
parfaitement explicables par le contexte macro-économique : *
\textbf{L'effet COVID-19 (2020) :} On observe une chute drastique du
nombre de sinistres déclarés (7 908 cas). Cela s'explique logiquement
par les confinements successifs et la réduction massive de la
circulation automobile cette année-là. * \textbf{Une reprise
post-pandémie :} Les années 2021 et 2022 retrouvent des niveaux de
sinistralité standards (autour de 9 500 à 11 300 sinistres). *
\textbf{Un pic exceptionnel en 2023 :} L'année 2023 affiche une
explosion avec plus de 15 000 sinistres. Cette hausse soudaine
nécessitera d'être croisée avec l'évolution de la taille de notre
portefeuille de contrats pour vérifier si elle est due à une acquisition
massive de nouveaux clients ou à une dégradation réelle du risque (ex:
événements climatiques).

\begin{Shaded}
\begin{Highlighting}[]
\CommentTok{\#1.Chargement de la table des contrats globaux }
\NormalTok{contrats\_globaux }\OtherTok{\textless{}{-}} \FunctionTok{readRDS}\NormalTok{(}\StringTok{"contrats\_exploration\_ok.rds"}\NormalTok{)}
\end{Highlighting}
\end{Shaded}

\begin{Shaded}
\begin{Highlighting}[]
\CommentTok{\#2. Préparation des données }
\NormalTok{data\_contrats\_annee }\OtherTok{\textless{}{-}}\NormalTok{ contrats\_globaux }\SpecialCharTok{\%\textgreater{}\%} 
  \FunctionTok{count}\NormalTok{(idxYear, }\AttributeTok{name =} \StringTok{"nombre\_contrats"}\NormalTok{)}

\FunctionTok{print}\NormalTok{(}\StringTok{"{-}{-}{-}NOMBRE DE CONTRAT PAR ANNEE{-}{-}{-}"}\NormalTok{)}
\end{Highlighting}
\end{Shaded}

\begin{verbatim}
## [1] "---NOMBRE DE CONTRAT PAR ANNEE---"
\end{verbatim}

\begin{Shaded}
\begin{Highlighting}[]
\FunctionTok{print}\NormalTok{(data\_contrats\_annee)}
\end{Highlighting}
\end{Shaded}

\begin{verbatim}
## # A tibble: 5 x 2
##   idxYear nombre_contrats
##     <dbl>           <int>
## 1    2019           55227
## 2    2020           56610
## 3    2021           58227
## 4    2022           61699
## 5    2023           69674
\end{verbatim}

\begin{Shaded}
\begin{Highlighting}[]
\CommentTok{\# 3. Création du graphique avec ggplot2}
\NormalTok{graph\_contrats }\OtherTok{\textless{}{-}} \FunctionTok{ggplot}\NormalTok{(data\_contrats\_annee, }\FunctionTok{aes}\NormalTok{(}\AttributeTok{x =} \FunctionTok{as.factor}\NormalTok{(idxYear), }\AttributeTok{y =}\NormalTok{ nombre\_contrats)) }\SpecialCharTok{+}
  \FunctionTok{geom\_bar}\NormalTok{(}\AttributeTok{stat =} \StringTok{"identity"}\NormalTok{, }\AttributeTok{fill =} \StringTok{"\#27ae60"}\NormalTok{) }\SpecialCharTok{+} \CommentTok{\# Un joli vert pour les contrats}
  \FunctionTok{geom\_text}\NormalTok{(}\FunctionTok{aes}\NormalTok{(}\AttributeTok{label =}\NormalTok{ nombre\_contrats), }\AttributeTok{vjust =} \SpecialCharTok{{-}}\FloatTok{0.5}\NormalTok{, }\AttributeTok{size =} \FloatTok{3.5}\NormalTok{, }\AttributeTok{color =} \StringTok{"black"}\NormalTok{) }\SpecialCharTok{+}
  \FunctionTok{labs}\NormalTok{(}
    \AttributeTok{title =} \StringTok{"Évolution de la taille du portefeuille"}\NormalTok{,}
    \AttributeTok{subtitle =} \StringTok{"Nombre de contrats par année d\textquotesingle{}exercice"}\NormalTok{,}
    \AttributeTok{x =} \StringTok{"Année d\textquotesingle{}exercice"}\NormalTok{,}
    \AttributeTok{y =} \StringTok{"Nombre de contrats"}
\NormalTok{  ) }\SpecialCharTok{+}
  \FunctionTok{theme\_minimal}\NormalTok{()}

\CommentTok{\# Affichage du graphique}
\FunctionTok{print}\NormalTok{(graph\_contrats)}
\end{Highlighting}
\end{Shaded}

\pandocbounded{\includegraphics[keepaspectratio]{03_Visualisations_files/figure-latex/unnamed-chunk-3-1.pdf}}

\subsection{Interprétation : Évolution du
portefeuille}\label{interpruxe9tation-uxe9volution-du-portefeuille}

La mise en perspective de la taille du portefeuille avec le nombre de
sinistres confirme nos hypothèses : * \textbf{Validation de l'effet
COVID-19 :} En 2020, le portefeuille de l'assureur a continué de croître
(56 610 contrats), ce qui prouve que la chute brutale des sinistres est
bien liée à une baisse de la circulation (confinements) et non à une
fuite des clients. * \textbf{Explication du pic de 2023 :} La forte
hausse des sinistres en 2023 s'explique en grande partie par une
croissance très importante du portefeuille, qui atteint près de 70 000
contrats. La volumétrie des risques assurés a donc mécaniquement fait
augmenter le nombre de déclarations.

\begin{Shaded}
\begin{Highlighting}[]
\CommentTok{\#1. Préparation des données : type de véhicule}
\NormalTok{data\_segment }\OtherTok{\textless{}{-}}\NormalTok{ contrats\_globaux }\SpecialCharTok{\%\textgreater{}\%} 
  \FunctionTok{count}\NormalTok{(vhSegment,}\AttributeTok{name =} \StringTok{"nombre\_vehicules"}\NormalTok{) }\SpecialCharTok{\%\textgreater{}\%}
  \FunctionTok{arrange}\NormalTok{(}\FunctionTok{desc}\NormalTok{(nombre\_vehicules))}

\FunctionTok{print}\NormalTok{(}\StringTok{"{-}{-}{-}DONNEES : REPARTITION PAR TYPE DE VEHICULE{-}{-}{-}"}\NormalTok{)}
\end{Highlighting}
\end{Shaded}

\begin{verbatim}
## [1] "---DONNEES : REPARTITION PAR TYPE DE VEHICULE---"
\end{verbatim}

\begin{Shaded}
\begin{Highlighting}[]
\FunctionTok{print}\NormalTok{(data\_segment)}
\end{Highlighting}
\end{Shaded}

\begin{verbatim}
## # A tibble: 4 x 2
##   vhSegment nombre_vehicules
##   <chr>                <int>
## 1 Citadine            120614
## 2 Familiale            75566
## 3 Compacte             59953
## 4 SUV                  45304
\end{verbatim}

\begin{Shaded}
\begin{Highlighting}[]
\CommentTok{\# Graphique Segment}
\NormalTok{graph\_segment }\OtherTok{\textless{}{-}} \FunctionTok{ggplot}\NormalTok{(data\_segment, }\FunctionTok{aes}\NormalTok{(}\AttributeTok{x =} \FunctionTok{reorder}\NormalTok{(vhSegment, }\SpecialCharTok{{-}}\NormalTok{nombre\_vehicules), }\AttributeTok{y =}\NormalTok{ nombre\_vehicules)) }\SpecialCharTok{+}
  \FunctionTok{geom\_bar}\NormalTok{(}\AttributeTok{stat =} \StringTok{"identity"}\NormalTok{, }\AttributeTok{fill =} \StringTok{"\#3498db"}\NormalTok{) }\SpecialCharTok{+} \CommentTok{\# Bleu clair}
  \FunctionTok{geom\_text}\NormalTok{(}\FunctionTok{aes}\NormalTok{(}\AttributeTok{label =}\NormalTok{ nombre\_vehicules), }\AttributeTok{vjust =} \SpecialCharTok{{-}}\FloatTok{0.5}\NormalTok{, }\AttributeTok{size =} \FloatTok{3.5}\NormalTok{) }\SpecialCharTok{+}
  \FunctionTok{labs}\NormalTok{(}\AttributeTok{title =} \StringTok{"Répartition des contrats par type de véhicule"}\NormalTok{, }\AttributeTok{x =} \StringTok{"Segment du véhicule"}\NormalTok{, }\AttributeTok{y =} \StringTok{"Nombre de contrats"}\NormalTok{) }\SpecialCharTok{+}
  \FunctionTok{theme\_minimal}\NormalTok{()}

\FunctionTok{print}\NormalTok{(graph\_segment)}
\end{Highlighting}
\end{Shaded}

\pandocbounded{\includegraphics[keepaspectratio]{03_Visualisations_files/figure-latex/unnamed-chunk-4-1.pdf}}

\begin{Shaded}
\begin{Highlighting}[]
\CommentTok{\#2.  Préparation de données : Alimentation}
\NormalTok{data\_energie }\OtherTok{\textless{}{-}}\NormalTok{ contrats\_globaux }\SpecialCharTok{\%\textgreater{}\%} 
  \FunctionTok{count}\NormalTok{(vhEnergy, }\AttributeTok{name =} \StringTok{"nombre\_vehicules"}\NormalTok{) }\SpecialCharTok{\%\textgreater{}\%}
  \FunctionTok{arrange}\NormalTok{(}\FunctionTok{desc}\NormalTok{(nombre\_vehicules))}
\FunctionTok{print}\NormalTok{(}\StringTok{"{-}{-}{-}DONNEE : REPARTITION PAR ENERGIE {-}{-}{-}"}\NormalTok{)}
\end{Highlighting}
\end{Shaded}

\begin{verbatim}
## [1] "---DONNEE : REPARTITION PAR ENERGIE ---"
\end{verbatim}

\begin{Shaded}
\begin{Highlighting}[]
\FunctionTok{print}\NormalTok{(data\_energie)}
\end{Highlighting}
\end{Shaded}

\begin{verbatim}
## # A tibble: 3 x 2
##   vhEnergy           nombre_vehicules
##   <chr>                         <int>
## 1 Essence                      135559
## 2 Diesel                       120352
## 3 Electrique/Hybride            45526
\end{verbatim}

\begin{Shaded}
\begin{Highlighting}[]
\CommentTok{\# Graphique Énergie}
\NormalTok{graph\_energie }\OtherTok{\textless{}{-}} \FunctionTok{ggplot}\NormalTok{(data\_energie, }\FunctionTok{aes}\NormalTok{(}\AttributeTok{x =} \FunctionTok{reorder}\NormalTok{(vhEnergy, }\SpecialCharTok{{-}}\NormalTok{nombre\_vehicules), }\AttributeTok{y =}\NormalTok{ nombre\_vehicules)) }\SpecialCharTok{+}
  \FunctionTok{geom\_bar}\NormalTok{(}\AttributeTok{stat =} \StringTok{"identity"}\NormalTok{, }\AttributeTok{fill =} \StringTok{"\#e67e22"}\NormalTok{) }\SpecialCharTok{+} \CommentTok{\# Orange}
  \FunctionTok{geom\_text}\NormalTok{(}\FunctionTok{aes}\NormalTok{(}\AttributeTok{label =}\NormalTok{ nombre\_vehicules), }\AttributeTok{vjust =} \SpecialCharTok{{-}}\FloatTok{0.5}\NormalTok{, }\AttributeTok{size =} \FloatTok{3.5}\NormalTok{) }\SpecialCharTok{+}
  \FunctionTok{labs}\NormalTok{(}\AttributeTok{title =} \StringTok{"Répartition des contrats selon l\textquotesingle{}alimentation"}\NormalTok{, }\AttributeTok{x =} \StringTok{"Type d\textquotesingle{}énergie"}\NormalTok{, }\AttributeTok{y =} \StringTok{"Nombre de contrats"}\NormalTok{) }\SpecialCharTok{+}
  \FunctionTok{theme\_minimal}\NormalTok{()}

\FunctionTok{print}\NormalTok{(graph\_energie)}
\end{Highlighting}
\end{Shaded}

\pandocbounded{\includegraphics[keepaspectratio]{03_Visualisations_files/figure-latex/unnamed-chunk-6-1.pdf}}

\subsection{Interprétation : Profil du parc
automobile}\label{interpruxe9tation-profil-du-parc-automobile}

L'analyse univariée des caractéristiques des véhicules met en évidence
le positionnement de notre portefeuille : * \textbf{Un parc
majoritairement urbain :} Les citadines représentent le segment le plus
important (120 614 contrats), suivies des familiales (75 566). Les SUV,
bien que populaires sur le marché, restent minoritaires chez
ENSAssuRances (45 304). * \textbf{Une motorisation thermique dominante
:} L'essence (135 559) et le diesel (120 352) équipent la très grande
majorité des véhicules. Néanmoins, on note la présence d'une part
significative de véhicules Électriques/Hybrides (plus de 45 000
contrats), témoignant du renouvellement progressif du parc vers des
mobilités plus vertes.

\begin{Shaded}
\begin{Highlighting}[]
\CommentTok{\#1. Préparation des données : roupe de véhicule (vhGroup)}

\NormalTok{data\_groupe }\OtherTok{\textless{}{-}}\NormalTok{ contrats\_globaux }\SpecialCharTok{\%\textgreater{}\%} 
  \FunctionTok{count}\NormalTok{(vhGroup, }\AttributeTok{name =} \StringTok{"nombre\_vehicules"}\NormalTok{) }\SpecialCharTok{\%\textgreater{}\%}
  \FunctionTok{arrange}\NormalTok{(}\FunctionTok{desc}\NormalTok{(nombre\_vehicules))}

\FunctionTok{print}\NormalTok{(}\StringTok{"{-}{-}{-} DONNÉES : TOP 10 DES GROUPES DE VÉHICULES {-}{-}{-}"}\NormalTok{)}
\end{Highlighting}
\end{Shaded}

\begin{verbatim}
## [1] "--- DONNÉES : TOP 10 DES GROUPES DE VÉHICULES ---"
\end{verbatim}

\begin{Shaded}
\begin{Highlighting}[]
\FunctionTok{print}\NormalTok{(}\FunctionTok{head}\NormalTok{(data\_groupe, }\DecValTok{10}\NormalTok{))}
\end{Highlighting}
\end{Shaded}

\begin{verbatim}
## # A tibble: 10 x 2
##    vhGroup nombre_vehicules
##      <dbl>            <int>
##  1      15            15514
##  2      13            15168
##  3      14            15146
##  4      12            15138
##  5      11            14852
##  6      16            14618
##  7      10            13864
##  8      17            13799
##  9       9            13691
## 10      18            13615
\end{verbatim}

\begin{Shaded}
\begin{Highlighting}[]
\CommentTok{\# Graphique : Distribution des groupes (On ne prend que le Top 15 pour que ce soit lisible)}
\NormalTok{graph\_groupe }\OtherTok{\textless{}{-}}\NormalTok{ data\_groupe }\SpecialCharTok{\%\textgreater{}\%}
  \FunctionTok{top\_n}\NormalTok{(}\DecValTok{15}\NormalTok{, nombre\_vehicules) }\SpecialCharTok{\%\textgreater{}\%}
  \FunctionTok{ggplot}\NormalTok{(}\FunctionTok{aes}\NormalTok{(}\AttributeTok{x =} \FunctionTok{reorder}\NormalTok{(}\FunctionTok{as.factor}\NormalTok{(vhGroup), }\SpecialCharTok{{-}}\NormalTok{nombre\_vehicules), }\AttributeTok{y =}\NormalTok{ nombre\_vehicules)) }\SpecialCharTok{+}
  \FunctionTok{geom\_bar}\NormalTok{(}\AttributeTok{stat =} \StringTok{"identity"}\NormalTok{, }\AttributeTok{fill =} \StringTok{"\#8e44ad"}\NormalTok{) }\SpecialCharTok{+} \CommentTok{\# Un joli violet}
  \FunctionTok{labs}\NormalTok{(}\AttributeTok{title =} \StringTok{"Distribution des véhicules par Groupe (Top 15)"}\NormalTok{, }\AttributeTok{x =} \StringTok{"Groupe du véhicule"}\NormalTok{, }\AttributeTok{y =} \StringTok{"Nombre de contrats"}\NormalTok{) }\SpecialCharTok{+}
  \FunctionTok{theme\_minimal}\NormalTok{()}

\FunctionTok{print}\NormalTok{(graph\_groupe)}
\end{Highlighting}
\end{Shaded}

\pandocbounded{\includegraphics[keepaspectratio]{03_Visualisations_files/figure-latex/unnamed-chunk-7-1.pdf}}

\begin{Shaded}
\begin{Highlighting}[]
\CommentTok{\#2. Préparation des données : Option Petit Rouleur}

\NormalTok{data\_km }\OtherTok{\textless{}{-}}\NormalTok{  contrats\_globaux }\SpecialCharTok{\%\textgreater{}\%}
  \FunctionTok{count}\NormalTok{(ctKM, }\AttributeTok{name =} \StringTok{"nombre\_contrats"}\NormalTok{) }\SpecialCharTok{\%\textgreater{}\%} \FunctionTok{arrange}\NormalTok{(}\FunctionTok{desc}\NormalTok{(nombre\_contrats))}

\FunctionTok{print}\NormalTok{(}\StringTok{"{-}{-}{-} DONNÉES : OPTION PETIT ROULEUR {-}{-}{-}"}\NormalTok{)}
\end{Highlighting}
\end{Shaded}

\begin{verbatim}
## [1] "--- DONNÉES : OPTION PETIT ROULEUR ---"
\end{verbatim}

\begin{Shaded}
\begin{Highlighting}[]
\FunctionTok{print}\NormalTok{(data\_km)}
\end{Highlighting}
\end{Shaded}

\begin{verbatim}
## # A tibble: 2 x 2
##   ctKM  nombre_contrats
##   <chr>           <int>
## 1 N              260973
## 2 O               40464
\end{verbatim}

\begin{Shaded}
\begin{Highlighting}[]
\CommentTok{\# Graphique : Option Petit Rouleur}
\NormalTok{graph\_km }\OtherTok{\textless{}{-}} \FunctionTok{ggplot}\NormalTok{(data\_km, }\FunctionTok{aes}\NormalTok{(}\AttributeTok{x =}\NormalTok{ ctKM, }\AttributeTok{y =}\NormalTok{ nombre\_contrats)) }\SpecialCharTok{+}
  \FunctionTok{geom\_bar}\NormalTok{(}\AttributeTok{stat =} \StringTok{"identity"}\NormalTok{, }\AttributeTok{fill =} \StringTok{"\#16a085"}\NormalTok{, }\AttributeTok{width =} \FloatTok{0.5}\NormalTok{) }\SpecialCharTok{+} \CommentTok{\# Vert émeraude}
  \FunctionTok{geom\_text}\NormalTok{(}\FunctionTok{aes}\NormalTok{(}\AttributeTok{label =}\NormalTok{ nombre\_contrats), }\AttributeTok{vjust =} \SpecialCharTok{{-}}\FloatTok{0.5}\NormalTok{, }\AttributeTok{size =} \DecValTok{4}\NormalTok{) }\SpecialCharTok{+}
  \FunctionTok{labs}\NormalTok{(}\AttributeTok{title =} \StringTok{"Répartition de l\textquotesingle{}option Petit Rouleur"}\NormalTok{, }\AttributeTok{x =} \StringTok{"Option Petit Rouleur"}\NormalTok{, }\AttributeTok{y =} \StringTok{"Nombre de contrats"}\NormalTok{) }\SpecialCharTok{+}
  \FunctionTok{theme\_minimal}\NormalTok{()}

\FunctionTok{print}\NormalTok{(graph\_km)}
\end{Highlighting}
\end{Shaded}

\pandocbounded{\includegraphics[keepaspectratio]{03_Visualisations_files/figure-latex/unnamed-chunk-9-1.pdf}}

\subsection{Interprétation : Caractéristiques techniques et
contractuelles}\label{interpruxe9tation-caractuxe9ristiques-techniques-et-contractuelles}

L'analyse des variables de tarification classiques révèle les points
suivants : * \textbf{Une homogénéité des véhicules assurés :} La
distribution des groupes de véhicules (\texttt{vhGroup}) montre une
forte concentration sur les segments intermédiaires (groupes 11 à 18).
ENSAssuRances assure donc majoritairement des véhicules grand public
standards, avec très peu de véhicules aux extrêmes (très bas de gamme ou
haut de gamme/sportifs). * \textbf{Une option kilométrique minoritaire
:} L'option ``Petit Rouleur'' (\texttt{ctKM}) n'est souscrite que par
environ 13,4 \% des assurés (40 464 contrats contre 260 973 pour l'usage
classique). Il sera intéressant d'analyser par la suite si cette
population présente effectivement une fréquence de sinistralité plus
faible, justifiant ainsi une tarification avantageuse.

\begin{Shaded}
\begin{Highlighting}[]
\CommentTok{\#1. chargement de la base sinistre }
\NormalTok{base\_complete }\OtherTok{\textless{}{-}} \FunctionTok{readRDS}\NormalTok{ (}\StringTok{"base\_complete\_clean.rds"}\NormalTok{)}
\end{Highlighting}
\end{Shaded}

\begin{Shaded}
\begin{Highlighting}[]
\CommentTok{\#{-}{-}{-}ANALYSE PAR ÂGE{-}{-}{-}}
\CommentTok{\# Préparation des données : on arrondit l\textquotesingle{}âge pour pouvoir les regrouper}
\NormalTok{data\_age }\OtherTok{\textless{}{-}}\NormalTok{ base\_complete }\SpecialCharTok{\%\textgreater{}\%}
  \FunctionTok{mutate}\NormalTok{(}
    \AttributeTok{age\_arrondi =} \FunctionTok{floor}\NormalTok{(drv1Age)}
\NormalTok{  ) }\SpecialCharTok{\%\textgreater{}\%} \FunctionTok{count}\NormalTok{(age\_arrondi,}\AttributeTok{name =} \StringTok{"nombre\_sinistres"}\NormalTok{)}

\FunctionTok{print}\NormalTok{(}\StringTok{"{-}{-}{-} DONNÉES : SINISTRES PAR ÂGE (Aperçu) {-}{-}{-}"}\NormalTok{)}
\end{Highlighting}
\end{Shaded}

\begin{verbatim}
## [1] "--- DONNÉES : SINISTRES PAR ÂGE (Aperçu) ---"
\end{verbatim}

\begin{Shaded}
\begin{Highlighting}[]
\FunctionTok{print}\NormalTok{(}\FunctionTok{head}\NormalTok{(data\_age, }\DecValTok{10}\NormalTok{))}
\end{Highlighting}
\end{Shaded}

\begin{verbatim}
## # A tibble: 10 x 2
##    age_arrondi nombre_sinistres
##          <dbl>            <int>
##  1          18               83
##  2          19              252
##  3          20              367
##  4          21              571
##  5          22              717
##  6          23              887
##  7          24              938
##  8          25             1097
##  9          26              971
## 10          27             1105
\end{verbatim}

\begin{Shaded}
\begin{Highlighting}[]
\CommentTok{\# Graphique : Courbe des sinistres par âge}
\NormalTok{graph\_age }\OtherTok{\textless{}{-}} \FunctionTok{ggplot}\NormalTok{(data\_age, }\FunctionTok{aes}\NormalTok{(}\AttributeTok{x =}\NormalTok{ age\_arrondi, }\AttributeTok{y =}\NormalTok{ nombre\_sinistres)) }\SpecialCharTok{+}
  \FunctionTok{geom\_line}\NormalTok{(}\AttributeTok{color =} \StringTok{"\#c0392b"}\NormalTok{, }\AttributeTok{size =} \DecValTok{1}\NormalTok{) }\SpecialCharTok{+} \CommentTok{\# Une belle ligne rouge}
  \FunctionTok{geom\_point}\NormalTok{(}\AttributeTok{color =} \StringTok{"\#c0392b"}\NormalTok{) }\SpecialCharTok{+}
  \FunctionTok{labs}\NormalTok{(}\AttributeTok{title =} \StringTok{"Nombre de sinistres en fonction de l\textquotesingle{}âge du conducteur"}\NormalTok{, }
       \AttributeTok{x =} \StringTok{"Âge du conducteur (années)"}\NormalTok{, }
       \AttributeTok{y =} \StringTok{"Nombre total de sinistres"}\NormalTok{) }\SpecialCharTok{+}
  \FunctionTok{theme\_minimal}\NormalTok{()}

\FunctionTok{print}\NormalTok{(graph\_age)}
\end{Highlighting}
\end{Shaded}

\pandocbounded{\includegraphics[keepaspectratio]{03_Visualisations_files/figure-latex/unnamed-chunk-11-1.pdf}}

\begin{Shaded}
\begin{Highlighting}[]
\CommentTok{\#{-}{-}{-} ANALYSE PAR SEXE}
\CommentTok{\# Préparation des données et calcul du pourcentage }
\NormalTok{data\_sexe }\OtherTok{\textless{}{-}}\NormalTok{ base\_complete }\SpecialCharTok{\%\textgreater{}\%} 
  \FunctionTok{count}\NormalTok{(drv1Sex,}\AttributeTok{name =}\StringTok{"nombre\_sinistres"}\NormalTok{) }\SpecialCharTok{\%\textgreater{}\%}
  \FunctionTok{mutate}\NormalTok{(}\AttributeTok{pourcentage =} \FunctionTok{round}\NormalTok{((nombre\_sinistres }\SpecialCharTok{/} \FunctionTok{sum}\NormalTok{(nombre\_sinistres)) }\SpecialCharTok{*} \DecValTok{100}\NormalTok{,}\DecValTok{1}\NormalTok{))}

\FunctionTok{print}\NormalTok{(}\StringTok{"{-}{-}{-} DONNÉES : SINISTRES PAR SEXE {-}{-}{-}"}\NormalTok{)}
\end{Highlighting}
\end{Shaded}

\begin{verbatim}
## [1] "--- DONNÉES : SINISTRES PAR SEXE ---"
\end{verbatim}

\begin{Shaded}
\begin{Highlighting}[]
\FunctionTok{print}\NormalTok{(data\_sexe)}
\end{Highlighting}
\end{Shaded}

\begin{verbatim}
## # A tibble: 2 x 3
##   drv1Sex nombre_sinistres pourcentage
##   <chr>              <int>       <dbl>
## 1 F                  23891        40.8
## 2 H                  34693        59.2
\end{verbatim}

\begin{Shaded}
\begin{Highlighting}[]
\CommentTok{\# Graphique : Répartition des sinistres par sexe}
\NormalTok{graph\_sexe }\OtherTok{\textless{}{-}} \FunctionTok{ggplot}\NormalTok{(data\_sexe, }\FunctionTok{aes}\NormalTok{(}\AttributeTok{x =}\NormalTok{ drv1Sex, }\AttributeTok{y =}\NormalTok{ nombre\_sinistres, }\AttributeTok{fill =}\NormalTok{ drv1Sex)) }\SpecialCharTok{+}
  \FunctionTok{geom\_bar}\NormalTok{(}\AttributeTok{stat =} \StringTok{"identity"}\NormalTok{, }\AttributeTok{width =} \FloatTok{0.5}\NormalTok{) }\SpecialCharTok{+}
  \FunctionTok{geom\_text}\NormalTok{(}\FunctionTok{aes}\NormalTok{(}\AttributeTok{label =} \FunctionTok{paste0}\NormalTok{(nombre\_sinistres, }\StringTok{" ("}\NormalTok{, pourcentage, }\StringTok{"\%)"}\NormalTok{)), }\AttributeTok{vjust =} \SpecialCharTok{{-}}\FloatTok{0.5}\NormalTok{, }\AttributeTok{size =} \DecValTok{4}\NormalTok{) }\SpecialCharTok{+}
  \FunctionTok{scale\_fill\_manual}\NormalTok{(}\AttributeTok{values =} \FunctionTok{c}\NormalTok{(}\StringTok{"F"} \OtherTok{=} \StringTok{"\#e74c3c"}\NormalTok{, }\StringTok{"H"} \OtherTok{=} \StringTok{"\#3498db"}\NormalTok{)) }\SpecialCharTok{+} \CommentTok{\# Couleurs distinctes}
  \FunctionTok{labs}\NormalTok{(}\AttributeTok{title =} \StringTok{"Nombre et pourcentage de sinistres selon le sexe"}\NormalTok{, }
       \AttributeTok{x =} \StringTok{"Sexe du conducteur"}\NormalTok{, }
       \AttributeTok{y =} \StringTok{"Nombre de sinistres"}\NormalTok{) }\SpecialCharTok{+}
  \FunctionTok{theme\_minimal}\NormalTok{() }\SpecialCharTok{+}
  \FunctionTok{theme}\NormalTok{(}\AttributeTok{legend.position =} \StringTok{"none"}\NormalTok{) }\CommentTok{\# On cache la légende inutile}

\FunctionTok{print}\NormalTok{(graph\_sexe)}
\end{Highlighting}
\end{Shaded}

\pandocbounded{\includegraphics[keepaspectratio]{03_Visualisations_files/figure-latex/unnamed-chunk-13-1.pdf}}

\subsection{Interprétation : Analyse du risque lié au
conducteur}\label{interpruxe9tation-analyse-du-risque-liuxe9-au-conducteur}

L'étude des sinistres selon le profil des assurés met en évidence des
marqueurs de risque classiques en assurance automobile : * \textbf{Le
risque ``Jeune Conducteur'' :} L'analyse de la courbe par âge montre une
augmentation très rapide du nombre d'accidents chez les jeunes (avec un
franchissement de la barre des 1 000 sinistres dès 25 ans). Ce volume
élevé confirme la nécessité d'une surprime pour les conducteurs novices.
* \textbf{La répartition par Sexe :} Les hommes déclarent la majorité
des sinistres (59,2 \%, soit 34 693 dossiers) contre 40,8 \% pour les
femmes. Toutefois, d'un point de vue actuariel, ce chiffre brut devra
être relativisé par la proportion d'hommes et de femmes dans le
portefeuille global (l'exposition) avant de pouvoir en tirer une
conclusion sur la sinistralité réelle (fréquence) par genre.

\begin{Shaded}
\begin{Highlighting}[]
\CommentTok{\#1. Analyse par antécédents de sinistres }
\NormalTok{data\_antecedents }\OtherTok{\textless{}{-}}\NormalTok{ base\_complete }\SpecialCharTok{\%\textgreater{}\%}
  \FunctionTok{count}\NormalTok{(claimsAnt,}\AttributeTok{name =} \StringTok{"nombre\_sinistres"}\NormalTok{)}

\FunctionTok{print}\NormalTok{(}\StringTok{"{-}{-}{-} DONNÉES : SINISTRES SELON LES ANTÉCÉDENTS {-}{-}{-}"}\NormalTok{)}
\end{Highlighting}
\end{Shaded}

\begin{verbatim}
## [1] "--- DONNÉES : SINISTRES SELON LES ANTÉCÉDENTS ---"
\end{verbatim}

\begin{Shaded}
\begin{Highlighting}[]
\FunctionTok{print}\NormalTok{(data\_antecedents)}
\end{Highlighting}
\end{Shaded}

\begin{verbatim}
## # A tibble: 3 x 2
##   claimsAnt nombre_sinistres
##       <dbl>            <int>
## 1         0            40969
## 2         1            11849
## 3         2             5766
\end{verbatim}

\begin{Shaded}
\begin{Highlighting}[]
\CommentTok{\# Graphique : Antécédents}
\NormalTok{graph\_antecedents }\OtherTok{\textless{}{-}} \FunctionTok{ggplot}\NormalTok{(data\_antecedents, }\FunctionTok{aes}\NormalTok{(}\AttributeTok{x =} \FunctionTok{as.factor}\NormalTok{(claimsAnt), }\AttributeTok{y =}\NormalTok{ nombre\_sinistres)) }\SpecialCharTok{+}
  \FunctionTok{geom\_bar}\NormalTok{(}\AttributeTok{stat =} \StringTok{"identity"}\NormalTok{, }\AttributeTok{fill =} \StringTok{"\#34495e"}\NormalTok{) }\SpecialCharTok{+} \CommentTok{\# Gris foncé très pro}
  \FunctionTok{geom\_text}\NormalTok{(}\FunctionTok{aes}\NormalTok{(}\AttributeTok{label =}\NormalTok{ nombre\_sinistres), }\AttributeTok{vjust =} \SpecialCharTok{{-}}\FloatTok{0.5}\NormalTok{, }\AttributeTok{size =} \DecValTok{4}\NormalTok{) }\SpecialCharTok{+}
  \FunctionTok{labs}\NormalTok{(}\AttributeTok{title =} \StringTok{"Nombre de sinistres selon les antécédents"}\NormalTok{,}
       \AttributeTok{x =} \StringTok{"Nombre de sinistres antérieurs (claimsAnt)"}\NormalTok{,}
       \AttributeTok{y =} \StringTok{"Total des sinistres actuels"}\NormalTok{) }\SpecialCharTok{+}
  \FunctionTok{theme\_minimal}\NormalTok{()}

\FunctionTok{print}\NormalTok{(graph\_antecedents)}
\end{Highlighting}
\end{Shaded}

\pandocbounded{\includegraphics[keepaspectratio]{03_Visualisations_files/figure-latex/unnamed-chunk-14-1.pdf}}

\begin{Shaded}
\begin{Highlighting}[]
\CommentTok{\#2. Analyse par segment commercial}
\NormalTok{data\_sin\_segment }\OtherTok{\textless{}{-}}\NormalTok{ base\_complete }\SpecialCharTok{\%\textgreater{}\%} 
  \FunctionTok{count}\NormalTok{(vhSegment, }\AttributeTok{name=}\StringTok{"nombre\_sinistres"}\NormalTok{) }\SpecialCharTok{\%\textgreater{}\%}
  \FunctionTok{arrange}\NormalTok{(}\FunctionTok{desc}\NormalTok{(nombre\_sinistres))}

\FunctionTok{print}\NormalTok{(}\StringTok{"{-}{-}{-} DONNÉES : SINISTRES PAR SEGMENT {-}{-}{-}"}\NormalTok{)}
\end{Highlighting}
\end{Shaded}

\begin{verbatim}
## [1] "--- DONNÉES : SINISTRES PAR SEGMENT ---"
\end{verbatim}

\begin{Shaded}
\begin{Highlighting}[]
\FunctionTok{print}\NormalTok{(data\_sin\_segment)}
\end{Highlighting}
\end{Shaded}

\begin{verbatim}
## # A tibble: 4 x 2
##   vhSegment nombre_sinistres
##   <chr>                <int>
## 1 Citadine             23645
## 2 Familiale            14431
## 3 Compacte             11731
## 4 SUV                   8777
\end{verbatim}

\begin{Shaded}
\begin{Highlighting}[]
\CommentTok{\# Graphique : Segment commercial}
\NormalTok{graph\_sin\_segment }\OtherTok{\textless{}{-}} \FunctionTok{ggplot}\NormalTok{(data\_sin\_segment, }\FunctionTok{aes}\NormalTok{(}\AttributeTok{x =} \FunctionTok{reorder}\NormalTok{(vhSegment, }\SpecialCharTok{{-}}\NormalTok{nombre\_sinistres), }\AttributeTok{y =}\NormalTok{ nombre\_sinistres)) }\SpecialCharTok{+}
  \FunctionTok{geom\_bar}\NormalTok{(}\AttributeTok{stat =} \StringTok{"identity"}\NormalTok{, }\AttributeTok{fill =} \StringTok{"\#d35400"}\NormalTok{) }\SpecialCharTok{+} \CommentTok{\# Orange sombre}
  \FunctionTok{geom\_text}\NormalTok{(}\FunctionTok{aes}\NormalTok{(}\AttributeTok{label =}\NormalTok{ nombre\_sinistres), }\AttributeTok{vjust =} \SpecialCharTok{{-}}\FloatTok{0.5}\NormalTok{, }\AttributeTok{size =} \DecValTok{4}\NormalTok{) }\SpecialCharTok{+}
  \FunctionTok{labs}\NormalTok{(}\AttributeTok{title =} \StringTok{"Nombre de sinistres par segment commercial"}\NormalTok{,}
       \AttributeTok{x =} \StringTok{"Segment du véhicule"}\NormalTok{,}
       \AttributeTok{y =} \StringTok{"Nombre de sinistres"}\NormalTok{) }\SpecialCharTok{+}
  \FunctionTok{theme\_minimal}\NormalTok{()}

\FunctionTok{print}\NormalTok{(graph\_sin\_segment)}
\end{Highlighting}
\end{Shaded}

\pandocbounded{\includegraphics[keepaspectratio]{03_Visualisations_files/figure-latex/unnamed-chunk-16-1.pdf}}

\subsection{Interprétation : Antécédents et Volume par
Segment}\label{interpruxe9tation-antuxe9cuxe9dents-et-volume-par-segment}

L'étude des volumes de sinistres nous apporte deux enseignements majeurs
: * \textbf{Le poids des ``bons profils'' :} La majorité absolue des
sinistres (40 969) est le fait d'assurés n'ayant aucun antécédent
récent. Le volume de ``récidivistes'' reste toutefois un point
d'attention pour la tarification. * \textbf{Le biais de volume sur les
véhicules :} Les citadines enregistrent le plus grand nombre de
sinistres bruts (23 645). Cependant, ce chiffre n'indique pas
nécessairement une plus grande dangerosité : il est le reflet direct de
l'omniprésence des citadines dans notre portefeuille (plus de 120 000
contrats). Pour évaluer le risque réel de chaque type de véhicule, il
est impératif de passer d'une logique de volume à une logique de
\textbf{fréquence} (ratio sinistres / contrats).

\begin{Shaded}
\begin{Highlighting}[]
\CommentTok{\#1. calcul du nombre total de contrats par Segment}

\NormalTok{expo\_segment }\OtherTok{\textless{}{-}}\NormalTok{ contrats\_globaux }\SpecialCharTok{\%\textgreater{}\%}
  \FunctionTok{count}\NormalTok{(vhSegment, }\AttributeTok{name =} \StringTok{"total\_contrats"}\NormalTok{)}

\CommentTok{\#2.calcul du nomre de sinistres par Segment}

\NormalTok{sin\_segment }\OtherTok{\textless{}{-}}\NormalTok{ base\_complete }\SpecialCharTok{\%\textgreater{}\%}
  \FunctionTok{count}\NormalTok{(vhSegment,}\AttributeTok{name =} \StringTok{"total\_sinistres"}\NormalTok{)}

\CommentTok{\# 3. Calcul de la FRÉQUENCE (Le vrai risque !)}
\NormalTok{risque\_segment }\OtherTok{\textless{}{-}}\NormalTok{ expo\_segment }\SpecialCharTok{\%\textgreater{}\%}
  \CommentTok{\# On fusionne les deux tableaux}
  \FunctionTok{inner\_join}\NormalTok{(sin\_segment, }\AttributeTok{by =} \StringTok{"vhSegment"}\NormalTok{) }\SpecialCharTok{\%\textgreater{}\%}
  \CommentTok{\# On calcule la fréquence en pourcentage : (Sinistres / Contrats) * 100}
  \FunctionTok{mutate}\NormalTok{(}\AttributeTok{frequence\_pct =} \FunctionTok{round}\NormalTok{((total\_sinistres }\SpecialCharTok{/}\NormalTok{ total\_contrats) }\SpecialCharTok{*} \DecValTok{100}\NormalTok{, }\DecValTok{2}\NormalTok{)) }\SpecialCharTok{\%\textgreater{}\%}
  \FunctionTok{arrange}\NormalTok{(}\FunctionTok{desc}\NormalTok{(frequence\_pct))}

\FunctionTok{print}\NormalTok{(}\StringTok{"{-}{-}{-} DONNÉES : FRÉQUENCE DE SINISTRALITÉ PAR SEGMENT (LE VRAI RISQUE) {-}{-}{-}"}\NormalTok{)}
\end{Highlighting}
\end{Shaded}

\begin{verbatim}
## [1] "--- DONNÉES : FRÉQUENCE DE SINISTRALITÉ PAR SEGMENT (LE VRAI RISQUE) ---"
\end{verbatim}

\begin{Shaded}
\begin{Highlighting}[]
\FunctionTok{print}\NormalTok{(risque\_segment)}
\end{Highlighting}
\end{Shaded}

\begin{verbatim}
## # A tibble: 4 x 4
##   vhSegment total_contrats total_sinistres frequence_pct
##   <chr>              <int>           <int>         <dbl>
## 1 Citadine          120614           23645          19.6
## 2 Compacte           59953           11731          19.6
## 3 SUV                45304            8777          19.4
## 4 Familiale          75566           14431          19.1
\end{verbatim}

\begin{Shaded}
\begin{Highlighting}[]
\CommentTok{\# 4. Graphique : Fréquence de sinistralité}
\NormalTok{graph\_risque }\OtherTok{\textless{}{-}} \FunctionTok{ggplot}\NormalTok{(risque\_segment, }\FunctionTok{aes}\NormalTok{(}\AttributeTok{x =} \FunctionTok{reorder}\NormalTok{(vhSegment, }\SpecialCharTok{{-}}\NormalTok{frequence\_pct), }\AttributeTok{y =}\NormalTok{ frequence\_pct)) }\SpecialCharTok{+}
  \FunctionTok{geom\_bar}\NormalTok{(}\AttributeTok{stat =} \StringTok{"identity"}\NormalTok{, }\AttributeTok{fill =} \StringTok{"\#c0392b"}\NormalTok{) }\SpecialCharTok{+} \CommentTok{\# Rouge pour indiquer le risque}
  \FunctionTok{geom\_text}\NormalTok{(}\FunctionTok{aes}\NormalTok{(}\AttributeTok{label =} \FunctionTok{paste0}\NormalTok{(frequence\_pct, }\StringTok{" \%"}\NormalTok{)), }\AttributeTok{vjust =} \SpecialCharTok{{-}}\FloatTok{0.5}\NormalTok{, }\AttributeTok{size =} \DecValTok{4}\NormalTok{, }\AttributeTok{fontface =} \StringTok{"bold"}\NormalTok{) }\SpecialCharTok{+}
  \FunctionTok{labs}\NormalTok{(}\AttributeTok{title =} \StringTok{"Identification des véhicules à risque"}\NormalTok{,}
       \AttributeTok{subtitle =} \StringTok{"Fréquence de sinistralité par segment (Sinistres / Contrats)"}\NormalTok{,}
       \AttributeTok{x =} \StringTok{"Segment du véhicule"}\NormalTok{,}
       \AttributeTok{y =} \StringTok{"Fréquence de sinistralité (\%)"}\NormalTok{) }\SpecialCharTok{+}
  \FunctionTok{theme\_minimal}\NormalTok{()}

\FunctionTok{print}\NormalTok{(graph\_risque)}
\end{Highlighting}
\end{Shaded}

\pandocbounded{\includegraphics[keepaspectratio]{03_Visualisations_files/figure-latex/unnamed-chunk-17-1.pdf}}

\subsection{Interprétation : Fréquence et Véhicules à
Risque}\label{interpruxe9tation-fruxe9quence-et-vuxe9hicules-uxe0-risque}

Le passage d'une analyse en volume à une analyse en fréquence révèle une
information capitale pour la tarification : * \textbf{L'illusion du
volume :} Bien que les Citadines génèrent de loin le plus grand nombre
de sinistres, leur fréquence de sinistralité (19,60 \%) est quasiment
identique à celle des autres segments (19,56 \% pour les Compactes,
19,09 \% pour les Familiales). * \textbf{Conclusion sur le risque :}
Dans le portefeuille d'ENSAssuRances, le segment commercial du véhicule
n'est pas un facteur discriminant majeur de la sinistralité. Le risque
est réparti de manière très homogène sur l'ensemble des catégories de
véhicules. La tarification devra donc s'appuyer sur d'autres facteurs
(comme l'âge ou les antécédents) plutôt que sur le seul type de
carrosserie.

\begin{Shaded}
\begin{Highlighting}[]
\FunctionTok{library}\NormalTok{(stringr) }\CommentTok{\# Pour manipuler les chaînes de caractères (le code INSEE)}

\CommentTok{\# 1. Extraction du département pour tous les contrats}
\NormalTok{contrats\_geo }\OtherTok{\textless{}{-}}\NormalTok{ contrats\_globaux }\SpecialCharTok{\%\textgreater{}\%}
  \CommentTok{\# On prend les 2 premiers caractères du code INSEE pour avoir le département}
  \FunctionTok{mutate}\NormalTok{(}\AttributeTok{departement =} \FunctionTok{str\_sub}\NormalTok{(ctINSEE, }\DecValTok{1}\NormalTok{, }\DecValTok{2}\NormalTok{)) }\SpecialCharTok{\%\textgreater{}\%}
  \CommentTok{\# On filtre pour ne garder que la France métropolitaine (01 à 95)}
  \FunctionTok{filter}\NormalTok{(departement }\SpecialCharTok{\%in\%} \FunctionTok{sprintf}\NormalTok{(}\StringTok{"\%02d"}\NormalTok{, }\DecValTok{1}\SpecialCharTok{:}\DecValTok{95}\NormalTok{))}

\CommentTok{\# 2. Calcul de l\textquotesingle{}exposition (nombre de contrats par département)}
\NormalTok{expo\_dept }\OtherTok{\textless{}{-}}\NormalTok{ contrats\_geo }\SpecialCharTok{\%\textgreater{}\%}
  \FunctionTok{count}\NormalTok{(departement, }\AttributeTok{name =} \StringTok{"total\_contrats"}\NormalTok{)}

\CommentTok{\# 3. Extraction du département pour les sinistres}
\NormalTok{sinistres\_geo }\OtherTok{\textless{}{-}}\NormalTok{ base\_complete }\SpecialCharTok{\%\textgreater{}\%}
  \FunctionTok{mutate}\NormalTok{(}\AttributeTok{departement =} \FunctionTok{str\_sub}\NormalTok{(ctINSEE, }\DecValTok{1}\NormalTok{, }\DecValTok{2}\NormalTok{)) }\SpecialCharTok{\%\textgreater{}\%}
  \FunctionTok{filter}\NormalTok{(departement }\SpecialCharTok{\%in\%} \FunctionTok{sprintf}\NormalTok{(}\StringTok{"\%02d"}\NormalTok{, }\DecValTok{1}\SpecialCharTok{:}\DecValTok{95}\NormalTok{)) }\SpecialCharTok{\%\textgreater{}\%}
  \FunctionTok{count}\NormalTok{(departement, }\AttributeTok{name =} \StringTok{"total\_sinistres"}\NormalTok{)}

\CommentTok{\# 4. Calcul de la Fréquence par Département}
\NormalTok{risque\_geo }\OtherTok{\textless{}{-}}\NormalTok{ expo\_dept }\SpecialCharTok{\%\textgreater{}\%}
  \FunctionTok{inner\_join}\NormalTok{(sinistres\_geo, }\AttributeTok{by =} \StringTok{"departement"}\NormalTok{) }\SpecialCharTok{\%\textgreater{}\%}
  \FunctionTok{mutate}\NormalTok{(}\AttributeTok{frequence\_pct =} \FunctionTok{round}\NormalTok{((total\_sinistres }\SpecialCharTok{/}\NormalTok{ total\_contrats) }\SpecialCharTok{*} \DecValTok{100}\NormalTok{, }\DecValTok{2}\NormalTok{)) }\SpecialCharTok{\%\textgreater{}\%}
  \CommentTok{\# On se concentre sur les départements où l\textquotesingle{}on a au moins 100 contrats pour que ce soit statistique}
  \FunctionTok{filter}\NormalTok{(total\_contrats }\SpecialCharTok{\textgreater{}} \DecValTok{100}\NormalTok{) }\SpecialCharTok{\%\textgreater{}\%}
  \FunctionTok{arrange}\NormalTok{(}\FunctionTok{desc}\NormalTok{(frequence\_pct))}

\CommentTok{\# Affichage du Top 10 des départements les plus risqués}
\FunctionTok{print}\NormalTok{(}\StringTok{"{-}{-}{-} TOP 10 DES DÉPARTEMENTS LES PLUS RISQUÉS (FRÉQUENCE EN \%) {-}{-}{-}"}\NormalTok{)}
\end{Highlighting}
\end{Shaded}

\begin{verbatim}
## [1] "--- TOP 10 DES DÉPARTEMENTS LES PLUS RISQUÉS (FRÉQUENCE EN %) ---"
\end{verbatim}

\begin{Shaded}
\begin{Highlighting}[]
\FunctionTok{print}\NormalTok{(}\FunctionTok{head}\NormalTok{(risque\_geo, }\DecValTok{10}\NormalTok{))}
\end{Highlighting}
\end{Shaded}

\begin{verbatim}
## # A tibble: 10 x 4
##    departement total_contrats total_sinistres frequence_pct
##    <chr>                <int>           <int>         <dbl>
##  1 45                    2526             597          23.6
##  2 69                    2720             609          22.4
##  3 77                    5016            1120          22.3
##  4 36                    2145             477          22.2
##  5 84                    2055             444          21.6
##  6 46                    1871             402          21.5
##  7 83                    2195             464          21.1
##  8 28                    3233             682          21.1
##  9 51                    5756            1207          21.0
## 10 90                    1463             306          20.9
\end{verbatim}

\subsection{Interprétation : Analyse Spatiale et Tarification
Géographique}\label{interpruxe9tation-analyse-spatiale-et-tarification-guxe9ographique}

L'analyse de la fréquence de sinistralité par département (code INSEE)
révèle des disparités territoriales majeures : * \textbf{Des zones de
sur-risque :} Le département 45 (Loiret) affiche la fréquence la plus
élevée du portefeuille (23,6 \%), suivi de près par les départements 69
(Rhône) et 77 (Seine-et-Marne). Ces fréquences dépassent largement la
moyenne globale du portefeuille (environ 19,3 \%). * \textbf{Impact
tarifaire :} La présence de départements très urbains ou traversés par
des axes routiers majeurs dans ce Top 10 confirme que la zone de
circulation est un facteur de risque déterminant. Cette cartographie
statistique justifie la mise en place d'un zonier tarifaire (surprime
pour les départements à forte fréquence) pour maintenir l'équilibre
technique d'ENSAssuRances.

\begin{Shaded}
\begin{Highlighting}[]
\FunctionTok{library}\NormalTok{(sf)}
\FunctionTok{library}\NormalTok{(ggplot2)}

\CommentTok{\# 1. CHARGEMENT DU SHAPEFILE}
\NormalTok{france\_sf }\OtherTok{\textless{}{-}} \FunctionTok{st\_read}\NormalTok{(}\StringTok{"ARRONDISSEMENT.shp"}\NormalTok{, }\AttributeTok{quiet =} \ConstantTok{TRUE}\NormalTok{) }\CommentTok{\# On ajoute quiet = TRUE pour éviter le texte de log dans le rapport final}

\CommentTok{\# 2. JOINTURE : Fusionner la carte avec nos statistiques}
\CommentTok{\# On utilise la bonne colonne CODE\_DEPT trouvée grâce à la fonction names()}
\NormalTok{carte\_data }\OtherTok{\textless{}{-}}\NormalTok{ france\_sf }\SpecialCharTok{\%\textgreater{}\%}
  \FunctionTok{left\_join}\NormalTok{(risque\_geo, }\AttributeTok{by =} \FunctionTok{c}\NormalTok{(}\StringTok{"CODE\_DEPT"} \OtherTok{=} \StringTok{"departement"}\NormalTok{))}

\CommentTok{\# 3. CRÉATION DE LA CARTE THÉMATIQUE}
\NormalTok{carte\_finale }\OtherTok{\textless{}{-}} \FunctionTok{ggplot}\NormalTok{(carte\_data) }\SpecialCharTok{+}
  \CommentTok{\# Tracé des polygones}
  \FunctionTok{geom\_sf}\NormalTok{(}\FunctionTok{aes}\NormalTok{(}\AttributeTok{fill =}\NormalTok{ frequence\_pct), }\AttributeTok{color =} \StringTok{"black"}\NormalTok{, }\AttributeTok{linewidth =} \FloatTok{0.1}\NormalTok{) }\SpecialCharTok{+}
  \CommentTok{\# Dégradé de couleurs personnalisé}
  \FunctionTok{scale\_fill\_gradient}\NormalTok{(}
    \AttributeTok{low =} \StringTok{"\#f1c40f"}\NormalTok{,      }\CommentTok{\# Jaune pour un risque faible}
    \AttributeTok{high =} \StringTok{"\#c0392b"}\NormalTok{,     }\CommentTok{\# Rouge foncé pour un risque élevé}
    \AttributeTok{na.value =} \StringTok{"\#ecf0f1"}\NormalTok{, }\CommentTok{\# Gris clair pour les zones sans données (DOM{-}TOM, etc.)}
    \AttributeTok{name =} \StringTok{"Fréquence (\%)"}
\NormalTok{  ) }\SpecialCharTok{+}
  \FunctionTok{labs}\NormalTok{(}
    \AttributeTok{title =} \StringTok{"Carte de la sinistralité en France"}\NormalTok{,}
    \AttributeTok{subtitle =} \StringTok{"Fréquence de sinistralité par département (Portefeuille ENSAssuRances)"}\NormalTok{,}
    \AttributeTok{caption =} \StringTok{"Données géographiques : IGN (GEOFLA)"}
\NormalTok{  ) }\SpecialCharTok{+}
  \CommentTok{\# Thème épuré pour une carte lisible}
  \FunctionTok{theme\_void}\NormalTok{() }\SpecialCharTok{+}
  \FunctionTok{theme}\NormalTok{(}\AttributeTok{legend.position =} \StringTok{"right"}\NormalTok{)}

\CommentTok{\# Affichage de la carte}
\FunctionTok{print}\NormalTok{(carte\_finale)}
\end{Highlighting}
\end{Shaded}

\pandocbounded{\includegraphics[keepaspectratio]{03_Visualisations_files/figure-latex/cartographie-risque-1.pdf}}

\end{document}
